\apendice{Especificación de Requisitos}

\section{Introducción}

Este documento define el comportamiento esperado de Batch Downloader antes de su implementación. Los requisitos se identifican para permitir trazabilidad con casos de uso, diseño y pruebas.

\section{Actores}

\begin{description}
	\item[Usuario anónimo:] consulta el catálogo, crea selecciones temporales y descarga bundles públicos dentro de cuotas.
	\item[Usuario registrado:] conserva bundles, los publica, asigna estrellas, solicita software y consulta sus trabajos.
	\item[Administrador:] mantiene catálogo, fuentes, dominios, resolvers y estados; revisa solicitudes y registros.
	\item[Sistema externo oficial:] página o servicio del editor desde el que se resuelve y descarga un instalador.
	\item[Proveedor semántico:] servicio de embeddings o LLM utilizado detrás de una interfaz controlada.
\end{description}

\section{Objetivos generales}

\begin{itemize}
	\item Reducir el trabajo necesario para descargar varias aplicaciones.
	\item Mantener el origen oficial y la trazabilidad de cada instalador.
	\item Evitar almacenamiento permanente y eliminar artefactos temporales.
	\item Permitir reutilizar selecciones mediante bundles.
	\item Mejorar el descubrimiento mediante búsqueda semántica segura.
	\item Proteger la infraestructura frente a URL maliciosas y uso abusivo.
\end{itemize}

\section{Catálogo de requisitos funcionales}

\begin{longtable}{p{0.09\textwidth}p{0.56\textwidth}p{0.13\textwidth}p{0.12\textwidth}}
\caption{Requisitos funcionales.}\label{tabla:requisitos-funcionales}\\
\toprule
\textbf{ID} & \textbf{Descripción} & \textbf{Prioridad} & \textbf{Actor}\\
\midrule
\endfirsthead
\toprule
\textbf{ID} & \textbf{Descripción} & \textbf{Prioridad} & \textbf{Actor}\\
\midrule
\endhead
RF-01 & Mostrar las aplicaciones activas con nombre, descripción, sistemas, arquitecturas, etiquetas y estado. & Alta & Todos\\
RF-02 & Buscar aplicaciones por coincidencia parcial de nombre. & Alta & Todos\\
RF-03 & Filtrar por sistema operativo, arquitectura y etiquetas, exigiendo todas o un mínimo de etiquetas. & Alta & Todos\\
RF-04 & Añadir o retirar aplicaciones de una selección temporal. & Alta & Todos\\
RF-05 & Mostrar tamaño estimado, fuentes y advertencias antes de generar el ZIP. & Media & Todos\\
RF-06 & Crear un trabajo asíncrono para generar el ZIP y devolver su identificador. & Alta & Todos\\
RF-07 & Consultar el progreso global y el estado individual de cada aplicación. & Alta & Todos\\
RF-08 & Permitir cancelar un trabajo que todavía no haya finalizado. & Media & Todos\\
RF-09 & Descargar un ZIP listo antes de su expiración. & Alta & Todos\\
RF-10 & Producir un ZIP parcial cuando alguna aplicación falle y exista al menos una descarga válida. & Alta & Sistema\\
RF-11 & Incluir manifiesto, documento informativo y errores individuales en el ZIP. & Alta & Sistema\\
RF-12 & Registrar y eliminar los archivos temporales al finalizar, cancelar o expirar. & Alta & Sistema\\
RF-13 & Registrar, iniciar y cerrar sesión con controles de seguridad. & Alta & Registrado\\
RF-14 & Crear, editar y eliminar bundles propios. & Alta & Registrado\\
RF-15 & Definir un bundle propio como privado o público y generar un enlace compartible si es público. & Alta & Registrado\\
RF-16 & Personalizar temporalmente cualquier bundle descargable sin modificar el original. & Alta & Todos\\
RF-17 & Clonar un bundle público como nuevo bundle del usuario. & Media & Registrado\\
RF-18 & Asignar o retirar una estrella de un bundle público ajeno, con un máximo de una por usuario. & Media & Registrado\\
RF-19 & Ordenar bundles por estrellas, fecha, nombre o etiquetas. & Media & Todos\\
RF-20 & Solicitar la incorporación de software indicando nombre y página oficial. & Media & Registrado\\
RF-21 & Añadir, editar, desactivar y reactivar aplicaciones, fuentes y etiquetas. & Alta & Administrador\\
RF-22 & Configurar dominios permitidos y estrategia de resolución por fuente. & Alta & Administrador\\
RF-23 & Probar un resolver y consultar sus redirecciones, metadatos y errores. & Alta & Administrador\\
RF-24 & Revisar trabajos, fallos de resolución, solicitudes y acciones administrativas. & Alta & Administrador\\
RF-25 & Ejecutar comprobaciones periódicas y marcar fuentes para revisión manual. & Media & Sistema\\
RF-26 & Buscar aplicaciones mediante lenguaje natural y devolver candidatos del catálogo. & Media & Todos\\
RF-27 & Proponer un bundle editable a partir de una consulta semántica. & Media & Todos\\
RF-28 & Explicar recomendaciones y resumir errores sin ejecutar acciones críticas. & Media & Todos\\
RF-29 & Enviar un aviso de finalización por correo cuando el usuario lo solicite y exista una dirección verificada. & Baja & Registrado\\
RF-30 & Permitir a un administrador crear y mantener bundles oficiales. & Media & Administrador\\
\bottomrule
\end{longtable}

\section{Requisitos no funcionales}

\begin{longtable}{p{0.09\textwidth}p{0.64\textwidth}p{0.17\textwidth}}
\caption{Requisitos no funcionales.}\label{tabla:requisitos-no-funcionales}\\
\toprule
\textbf{ID} & \textbf{Descripción verificable} & \textbf{Categoría}\\
\midrule
\endfirsthead
\toprule
\textbf{ID} & \textbf{Descripción verificable} & \textbf{Categoría}\\
\midrule
\endhead
RNF-01 & Solo se aceptarán descargas HTTPS desde dominios permitidos; se bloquearán redes privadas, locales y reservadas antes y después de redirecciones. & Seguridad\\
RNF-02 & Ningún instalador se almacenará permanentemente. Los temporales y ZIP expirados se eliminarán mediante políticas automáticas verificables. & Privacidad\\
RNF-03 & Los nombres incorporados al ZIP no podrán contener rutas, secuencias de escape, nombres reservados ni colisiones sin resolver. & Seguridad\\
RNF-04 & Se aplicarán límites configurables de tamaño, aplicaciones, tiempo, redirecciones, reintentos y concurrencia. & Seguridad\\
RNF-05 & Los permisos se comprobarán en el backend para cada operación protegida. & Seguridad\\
RNF-06 & Las contraseñas se almacenarán con un algoritmo de derivación resistente y nunca aparecerán en registros. & Seguridad\\
RNF-07 & Las operaciones administrativas y el uso de claves de API quedarán auditados. & Trazabilidad\\
RNF-08 & La API se documentará mediante OpenAPI y mantendrá errores estructurados. & Mantenibilidad\\
RNF-09 & La creación de un trabajo responderá sin esperar a que terminen las descargas. & Rendimiento\\
RNF-10 & La interfaz mostrará progreso y errores con lenguaje comprensible y será operable con teclado. & Usabilidad\\
RNF-11 & Los componentes se ejecutarán en contenedores y podrán configurarse mediante variables y secretos externos. & Portabilidad\\
RNF-12 & La API y los workers podrán replicarse de forma independiente. & Escalabilidad\\
RNF-13 & El sistema dispondrá de métricas de trabajos, colas, errores, temporales y tiempos por fase. & Observabilidad\\
RNF-14 & El LLM solo recibirá el contexto mínimo, su salida se validará y el sistema continuará sin él cuando falle. & Fiabilidad\\
RNF-15 & Las pruebas automatizadas cubrirán reglas críticas de URL, dominios, permisos, ZIP y estados. & Calidad\\
RNF-16 & Los datos personales se conservarán durante plazos definidos y podrán ser exportados o eliminados conforme a la normativa. & Legal\\
RNF-17 & Las páginas principales deberán adaptarse a escritorio y dispositivos móviles sin ocultar información esencial. & Usabilidad\\
RNF-18 & La interfaz distinguirá claramente página oficial, dominio final, tamaño conocido o estimado y fecha de comprobación. & Transparencia\\
\bottomrule
\end{longtable}

\section{Reglas de negocio}

\begin{enumerate}
	\item Una selección de dos o más aplicaciones se tratará como bundle temporal, aunque no se guarde.
	\item Un usuario solo podrá mantener una estrella activa por bundle.
	\item La personalización de un bundle ajeno no alterará el original.
	\item Solo el propietario podrá editar o eliminar permanentemente su bundle; un administrador podrá moderarlo según las condiciones de uso.
	\item Un bundle oficial solo podrá ser creado o marcado como tal por un administrador.
	\item Un bundle privado no será accesible mediante su antiguo enlace público.
	\item La fuente de verdad del catálogo será MySQL; pgvector mantendrá una representación derivada.
	\item Una sugerencia del LLM nunca evitará la validación normal de aplicaciones y descargas.
	\item El trabajo parcial será válido cuando incluya al menos un instalador y describa todos los fallos.
	\item Un trabajo listo dejará de estar disponible cuando se descargue, se elimine o venza su plazo.
\end{enumerate}

\section{Especificación de casos de uso}

\subsection{CU-01 Generar ZIP}

\begin{table}[H]
\centering
\begin{tabularx}{\textwidth}{p{0.23\textwidth}X}
\toprule
\textbf{Actor} & Usuario anónimo o registrado\\
\textbf{Requisitos} & RF-01 a RF-12; RNF-01 a RNF-05\\
\textbf{Precondición} & Existe al menos una aplicación seleccionada y no se superan las cuotas iniciales.\\
\textbf{Flujo principal} & 1. El usuario revisa selección, sistema y arquitectura.\newline
2. El sistema estima tamaño y muestra fuentes.\newline
3. El usuario confirma.\newline
4. La API valida y crea el trabajo.\newline
5. El worker resuelve y descarga cada aplicación.\newline
6. El sistema genera el ZIP y notifica que está listo.\newline
7. El usuario descarga el archivo.\\
\textbf{Alternativas} & Una aplicación puede fallar y producirse un ZIP parcial. El usuario puede cancelar. Si todas fallan, el trabajo termina como fallido.\\
\textbf{Postcondición} & El resultado queda disponible temporalmente y los artefactos se eliminan al expirar.\\
\bottomrule
\end{tabularx}
\caption{CU-01 Generar ZIP.}
\label{tabla:cu-generar-zip}
\end{table}

\subsection{CU-02 Crear y publicar bundle}

\begin{table}[H]
\centering
\begin{tabularx}{\textwidth}{p{0.23\textwidth}X}
\toprule
\textbf{Actor} & Usuario registrado\\
\textbf{Requisitos} & RF-14, RF-15 y RF-30\\
\textbf{Precondición} & El usuario ha iniciado sesión y existe una selección válida.\\
\textbf{Flujo principal} & 1. Introduce nombre y descripción.\newline
2. Revisa las aplicaciones.\newline
3. Guarda el bundle como privado.\newline
4. Opcionalmente cambia la visibilidad a pública.\newline
5. El sistema genera un enlace compartible.\\
\textbf{Excepciones} & Nombre inválido, aplicaciones eliminadas o intento de marcarlo como oficial sin permisos.\\
\textbf{Postcondición} & El bundle queda asociado a su propietario y conserva una composición persistente.\\
\bottomrule
\end{tabularx}
\caption{CU-02 Crear y publicar bundle.}
\label{tabla:cu-crear-bundle}
\end{table}

\subsection{CU-03 Personalizar un bundle público}

\begin{table}[H]
\centering
\begin{tabularx}{\textwidth}{p{0.23\textwidth}X}
\toprule
\textbf{Actor} & Cualquier usuario\\
\textbf{Requisitos} & RF-16 y RF-17\\
\textbf{Precondición} & El bundle es público o pertenece al usuario.\\
\textbf{Flujo principal} & 1. El usuario abre el bundle.\newline
2. Añade o retira aplicaciones para esa descarga.\newline
3. Genera el ZIP con la composición temporal.\\
\textbf{Postcondición} & El bundle original no cambia. Un usuario registrado puede elegir clonar la composición como uno nuevo.\\
\bottomrule
\end{tabularx}
\caption{CU-03 Personalizar un bundle público.}
\label{tabla:cu-personalizar-bundle}
\end{table}

\subsection{CU-04 Asignar estrella}

\begin{table}[H]
\centering
\begin{tabularx}{\textwidth}{p{0.23\textwidth}X}
\toprule
\textbf{Actor} & Usuario registrado\\
\textbf{Requisitos} & RF-18 y RF-19\\
\textbf{Precondición} & El bundle es público y el usuario puede valorarlo.\\
\textbf{Flujo principal} & El usuario activa o retira la estrella. El backend crea o elimina la relación única y actualiza el contador.\\
\textbf{Excepciones} & Bundle privado, inexistente o restricción de valorar el propio bundle.\\
\textbf{Postcondición} & El ranking refleja el nuevo número de estrellas.\\
\bottomrule
\end{tabularx}
\caption{CU-04 Asignar estrella.}
\label{tabla:cu-estrella}
\end{table}

\subsection{CU-05 Mantener una fuente}

\begin{table}[H]
\centering
\begin{tabularx}{\textwidth}{p{0.23\textwidth}X}
\toprule
\textbf{Actor} & Administrador\\
\textbf{Requisitos} & RF-21 a RF-25\\
\textbf{Precondición} & El administrador está autenticado.\\
\textbf{Flujo principal} & 1. Selecciona una aplicación.\newline
2. Configura sistema, arquitectura, página oficial, dominios y resolver.\newline
3. Ejecuta una prueba.\newline
4. Revisa URL final, redirecciones y metadatos.\newline
5. Guarda la fuente si el resultado es válido.\\
\textbf{Excepciones} & Dominio no permitido, red privada, contenido inesperado o resolución fallida.\\
\textbf{Postcondición} & La configuración queda versionada y la acción se registra.\\
\bottomrule
\end{tabularx}
\caption{CU-05 Mantener una fuente.}
\label{tabla:cu-mantener-fuente}
\end{table}

\subsection{CU-06 Buscar mediante lenguaje natural}

\begin{table}[H]
\centering
\begin{tabularx}{\textwidth}{p{0.23\textwidth}X}
\toprule
\textbf{Actor} & Cualquier usuario\\
\textbf{Requisitos} & RF-26 a RF-28; RNF-14\\
\textbf{Precondición} & La búsqueda semántica está habilitada.\\
\textbf{Flujo principal} & 1. El usuario escribe una necesidad.\newline
2. El sistema genera un embedding y recupera candidatos.\newline
3. Aplica filtros de estado y compatibilidad.\newline
4. Opcionalmente solicita una explicación al LLM.\newline
5. Valida y muestra los resultados.\\
\textbf{Alternativas} & Si falla el LLM, se muestran los candidatos sin explicación. Si no hay candidatos, se informa sin inventar aplicaciones.\\
\textbf{Postcondición} & El usuario puede añadir resultados a su selección o generar un borrador de bundle.\\
\bottomrule
\end{tabularx}
\caption{CU-06 Buscar mediante lenguaje natural.}
\label{tabla:cu-busqueda-semantica}
\end{table}

\section{Trazabilidad inicial}

\begin{table}[H]
\centering
\small
\begin{tabularx}{\textwidth}{p{0.23\textwidth}p{0.31\textwidth}X}
\toprule
\textbf{Área} & \textbf{Requisitos} & \textbf{Pruebas previstas}\\
\midrule
Catálogo y selección & RF-01 a RF-05 & API, filtros, compatibilidad y pruebas E2E.\\
Descarga y ZIP & RF-06 a RF-12 & Integración, descargas simuladas, ZIP parcial, cancelación y limpieza.\\
Usuarios y bundles & RF-13 a RF-19, RF-30 & Autorización, propiedad, enlace, clonación y estrella única.\\
Administración & RF-20 a RF-25 & Permisos, auditoría, resolver controlado y estados.\\
Semántica & RF-26 a RF-28 & Recuperación, filtros, esquema inválido y caída del proveedor.\\
Calidad transversal & RNF-01 a RNF-18 & Seguridad, carga, accesibilidad, observabilidad y privacidad.\\
\bottomrule
\end{tabularx}
\caption{Trazabilidad entre áreas y pruebas.}
\label{tabla:trazabilidad}
\end{table}
